\Qhead{4.8}{Генерация волн в СВЧ-диапазоне}{Генерация волн в СВЧ-диапазоне. Принцип работы и устройство лампы бегущей и обратной волны, магнетрона и клистрона. Отрицательное дифференциальное сопротивление и генераторы СВЧ на полевых транзисторах, туннельных диодах, диодах Ганна и лавиннопролетных диодах. Эффект Джозефсона.}
\Subfile{4.8.1}{sections/04_generation/q08/s01}
\Subfile{4.8.2}{sections/04_generation/q08/s02}
\Subfile{4.8.3}{sections/04_generation/q08/s03}

\recap{Вакуумные приборы используют группировку электронов и фазовый энергообмен; твердотельные — активный динамический отклик и обратную связь. Джозефсоновский переход напрямую связывает напряжение со скоростью фазы.}
\board{Клистронные сгустки; противоположные направления энергии и пучка в ЛОВ; отрицательная проводимость; две формулы Джозефсона.}
\oral{Почему магнитное поле магнетрона необходимо, хотя само не совершает работы?}
\bait{Можно оставить различие статической и динамической отрицательной проводимости: транзисторный генератор не требует падающей статической ВАХ.}
\sources{Carter~\cite{carter}, разделы о клистроне и магнетроне; практикумная запись ЛОВ~\cite{bwo}: принцип перестройки, без заимствования спорных степенных зависимостей; Биленко и соавт.~\cite{bilenko}, активные элементы. Джозефсоновские соотношения сопоставлены с NIST~\cite{nistjosephson}.}
